% A Log-Tree Top-K Engine for a NoC Mesh -- owned by briefs/04-system.md.
% Evidence: C2-1..C2-5, C3-1..C3-6 (evidence.md sections 1.3, 1.4), H11;
% forecast.md sections 1-5; rowskip-implementation.md sections 1-7;
% contrast rows from related-work.md section 3.
% Anonymity: kernels described structurally; no repo paths, no hashes.
\section{A Log-Tree Top-K Engine for a NoC Mesh}
\label{sec:system}

The map of Section~\ref{sec:design-space} leaves two live moves. First,
parallelize the data-oblivious corner (A)+(C)+(E): split the row across
cores and pay a log-depth rendezvous instead of a shared-memory partition.
Second, import the shrinking axis (D) in the only form that survives
without scatter: skip work, never move data. This section builds both:
Sections~\ref{sec:tree}--\ref{sec:costmodel} the column-parallel log-tree
operator (\topkop), Sections~\ref{sec:skip-mech}--\ref{sec:forecast}
chunk-skip and the calibrated forecast that eliminated its losing variant
pre-build.

\subsection{Operator Structure}
\label{sec:tree}

Figure~\ref{fig:operator} shows the operator on one row of $N$ elements,
using the notation of Section~\ref{sec:background}: the leaf window $M$
snaps to the LLK sort width (\S\ref{sec:bg-incumbent}), and the row
splits into
$C = N/M$ chunks. The program places $P$ slices on a rectangle of cores;
each slice streams $\lceil C/P \rceil$ chunks through a bitonic leaf that
sorts the arriving chunk, merges it into a resident $M$-wide descending
window, and rebuilds the window order. A $\lceil \log_2 P \rceil$-level
in-place merge tree --- the mesh-shaped sibling of FPGA merge-tree
sorters~\cite{samardzic2020bonsai}, applied to selection --- then folds
the $P$ windows to a root core, which emits the top-$K$ values and
indices; every level halves the live cores and merges two $M$-wide
sequences.

The synchronization fabric is deliberately austere. Tree levels map
one-to-one onto NoC semaphore levels; each core owns a single receive
circular buffer sized for two $M$-wide sequences; the writer kernels
extend to seven partner slots, enough for $P{=}128$. The operator uses no atomics, no
global barrier, and no dedicated DRAM scratch
(\S\ref{sec:bg-problem})~\cite{alabi2012kselection,li2024radik}. The tree
replaced a serial merge-chain ancestor that folded slices one at a time:
at fixed core counts the swap cut $k{=}2048$, $N{=}65{,}536$ from
69.6\us{} to 41.9\us{} (1.66$\times$) and $k{=}512$, $N{=}262{,}144$ from
64.2\us{} to 32.0\us{} (2.01$\times$), and it made runtime monotone in
$P$ --- the property the cost model below then exploits.

\begin{figure}[t]
  \centering
  \includegraphics[width=0.75\linewidth]{fig-s1-operator.pdf}
  \caption{Column-parallel log-tree top-$k$ operator.}
  \label{fig:operator}
\end{figure}

\subsection{Cost Model and Rectangle Embedding}
\label{sec:costmodel}

The factory chooses $P$ by minimizing a two-term model in leaf merge
units:
\begin{equation}
  \mathrm{cost}(P) \;=\; 2\,\big\lceil C/P \big\rceil
  \;+\; \big\lceil \log_2 P \big\rceil,
  \qquad P \in [2,\, \min(C, 128)],
  \label{eq:costmodel}
\end{equation}
with ties resolved toward fewer cores. The factor of two prices the leaf
honestly: every chunk after the first costs $\approx$2 merge units (sort,
merge, rebuild) while chunk~0 costs one. The unit constant is measured,
not assumed: 1.46\us{} per merge unit at $M{=}512$ and 5.51\us{} at
$M{=}2048$ ($\approx$1{,}966 and $\approx$7{,}440
cycles\footnote{\label{fn:clock}Cycle forms of Tracy-measured times assume
a 1.35\,GHz busy clock; the clock under load was not captured (idle
800\,MHz measured). Section~\ref{sec:evaluation} states the policy; all
cycle$\rightarrow$time conversions in this section carry it.}). The leaf
term is mesh-selection theory's multi-packet
regime~\cite{krizanc1992optimal,krizanc1993fast,multipacket1992} with
measured constants; the log term is the mesh's rendezvous price.

Minimizing Eq.~\ref{eq:costmodel} is only half the placement problem; the
winner must also embed in the 13$\times$10 worker grid as a contiguous
$a{\times}b$ rectangle. An earlier full-rows-only fit silently clamped
$P$: at $k{=}2048$, $N{=}65{,}536$ it settled for 13$\times$2 = 26 cores
while an 8$\times$4 = 32-core rectangle sat free. Searching all
cost-optimal rectangles instead wins 8--27\% at the op level across the
re-pinned cells: 41.9$\rightarrow$34.4\us{} (26$\rightarrow$32 cores),
58.6$\rightarrow$49.6\us{} (52$\rightarrow$64), 32.0$\rightarrow$23.3\us{}
(52$\rightarrow$104), and 15.0$\rightarrow$13.1\us{}
(52$\rightarrow$64). On a harvested mesh the
embedding search is part of the cost model, not a detail below it.
Section~\ref{sec:eval-scaling} defends the model against the measured
15-point $P$-sweep, including the ceiling-term bump it predicts.

\subsection{Chunk-Skip: Mechanism and Soundness}
\label{sec:skip-mech}

Ingredient (D) --- shrinking the touched data --- survives on this
machine only if nothing has to move. Chunk-skip is that import: a
per-chunk predicate that decides whether the leaf pipeline needs to run
at all. Let $T$ be the running $K$-th largest survivor, read from the
resident sorted window. The rule is strict:
skip chunk $c$ iff $\max(\text{chunk}_c) < T$. Chunk~0 always seeds the
window; $T$ is monotone nondecreasing because merges only improve the
window; boundary ties ($\max = T$) are never skipped; and the decision is
a pure function of the input, so the output is deterministic per input.

Soundness admits a short argument: at most $K{-}1$ elements ever exceed
the row's final $K$-th largest value $v_K$, so $T \le v_K$ at all times;
any member $x$ of an exact top-$K$ set satisfies $x \ge v_K \ge T$, so
the strict test never fires on $x$'s chunk, and $x$, sitting within the
top-$K \le$ top-$M$, survives every $M$-wide merge to the root. The
released kernel source carries the full 288-line proof as a header
comment. The contrast with the GPU running-threshold family is
threefold: WarpSelect thresholds per element, with no analytic skip
law~\cite{johnson2019billion}; GVR guesses its threshold from the
previous decode step and must count, verify, and refine after a
misprediction~\cite{gvr2026}; chunk-skip prunes at $M$-element
granularity, sound by construction --- no verification pass exists to
pay for --- with a distribution-free law and compile-time gate
(\S\ref{sec:skip-law}) and a pre-build forecast (\S\ref{sec:forecast})
that neither GPU mechanism has. Adversarial validation
backs the proof: a 36-check battery (all-equal, ascending, last-chunk
winners, boundary ties replicated across 64 chunks) is bit-exact in both
gated and ungated configurations, and a 20-launch hang battery passed
twice.

\subsection{The Skip Law and Its Compile-Time Gate}
\label{sec:skip-law}

How often the rule fires is a rank statistic, not a distribution
property. For i.i.d.\ inputs, chunk $c{+}1$ is skippable iff the top-$K$
of the first $(c{+}1)M$ elements all lie in the first $cM$:
\begin{equation}
  \Pr[\text{skip at } c] \;=\;
  \binom{cM}{K} \Big/ \binom{(c+1)M}{K}
  \;\approx\; e^{-K/(c+1)},
  \label{eq:skiplaw}
\end{equation}
identical for normal, uniform, and softmax-shaped inputs because only
ranks enter. The law is unforgiving at short streams: a
conservative $K{=}512$ threshold skips with probability
$2\times10^{-307}$ at $c{=}1$ and reaches only 1.8\% at $c{=}128$, while
an aggressive user-$K{=}32$ threshold --- sound by the argument of
Section~\ref{sec:skip-mech} even though the resident window is $M$-wide
--- reaches 37\% at $c{=}32$ and 78\% at $c{=}128$.

The law hands us the gate. Below $c = K/4$ the skip probability is under
$e^{-4} \approx 1.8\%$, so the compile-time gate sets the first tested
chunk to $\max(2, K/4)$ --- a derivation from Eq.~\ref{eq:skiplaw}, not a
tuned constant (the floor of two: the threshold address needs the
post-rebuild window layout of Section~\ref{sec:characterization}, which
exists only after the first merge). The gate zeroes the measured ungated overhead ($+6.8\%$ at
$K{=}512$ over 128 chunks, $+1.8\%$ at $K{=}1536$ over 25 chunks) while
forfeiting under one expected skip per row. The decision itself is cheap:
$\approx$150\,ns per tested chunk at $M{=}512$ and $\approx$350\,ns at
$M{=}2048$, derived from the ungated A/B (126 tested chunks cost
$+18.9$\us) and consistent with the 81-cycle rendezvous plus fold
components of Section~\ref{sec:design-space}
(footnote~\ref{fn:clock}) --- against $\approx$2.0\us{} and
$\approx$5.4\us{} saved per skipped chunk.

\subsection{Forecast Before Build: The Variant We Did Not Build}
\label{sec:forecast}

Before any kernel was written, an exact host-side simulation --- bf16
sign-magnitude streaming, stimulus matched to the sweep harness,
calibrated to reproduce all five pinned measured cells to
$\pm 0.0$\us{} by construction --- priced chunk-skip inside the
column-parallel tree. The verdict was a no-go with the weight of a
result: 0.00\% skip on every realistic distribution on every cell, a
predicted $+0.1$--$0.4$\us{} pure regression from the test overhead. The
cause is structural: \emph{the rectangle split and the skip law fight
over the same resource, stream length $c$} --- the cost-minimizing split
caps each slice at $\lceil C/P \rceil = 1$--$5$ chunks while
Eq.~\ref{eq:skiplaw} needs $c \gtrsim K/4$. The split already won.

The build therefore went where the law points: the row-parallel path, in
which each core streams all $C$ chunks of its own row --- long streams,
small $K$, exactly the shapes the stock small-$k$ routing ships to this
path. The forecast's redirect predicted 54.3\% skip at $K{=}32$ on a
128-chunk stream. At rows${=}2$, $K{=}32$, $N{=}65{,}536$ the measured
time falls 1.82$\times$ ($-45.1\%$), and at rows${=}8$ by 1.57$\times$
($-36.3\%$) --- both on 130-core programs whose active cores number
$\min(\text{rows}, 130)$, so the rows${=}2$ cell exercises two cores, not
the grid. Cells where the law forbids a win are protected by the gate, and the
column-parallel guard cell holds with byte-identical tree kernels
(Section~\ref{sec:eval-skip} shows the A/B bars, quotes the guard
numbers, and carries the epistemic caveat). The forecast said no;
we did not build that variant --- Section~\ref{sec:methodology} returns to
forecast-before-build as methodology.
